\chapter*{序言：先做出一个，再谈智能}
\addcontentsline{toc}{chapter}{序言：先做出一个，再谈智能}

过去几年里，“Agent”迅速从论文术语变成产品标签。聊天机器人加一个搜索按钮可以叫 Agent，能连续工作数小时的软件工程系统也叫 Agent；固定的自动化流程、多角色对话、浏览器操作、个人助理，常常被塞进同一个词里。结果是：资源越来越多，真正能把系统做可靠的人却没有同比增加。

这本书试图解决一个朴素问题：\term{怎样把一个模糊的 Agent 想法，变成可运行、可验证、可维护的系统？}

答案不会从框架排行榜开始。我们先从最小循环开始：模型观察环境，选择动作，调用工具，读取真实反馈，再决定下一步。随后逐层加入工具契约、状态、上下文、记忆、检索、规划、评测、权限和可观测性。每增加一层复杂度，都要回答两个问题：它解决了什么已被观察到的问题？我们如何证明它确实有用？

\begin{takeaway}[这本书的四条主线]
\begin{enumerate}
  \item \textbf{任务主线}：从任务、用户和验收标准出发，不从模型或框架出发。
  \item \textbf{反馈主线}：让 Agent 从工具、测试和真实环境得到可验证反馈。
  \item \textbf{可靠性主线}：失败恢复、权限、评测和追踪与功能同时设计。
  \item \textbf{作品主线}：每一部分都留下可以运行、展示和复盘的产物。
\end{enumerate}
\end{takeaway}

\section*{这本书适合谁}

\begin{itemize}
  \item \textbf{开发者}：会写基础 Python，希望系统掌握 Agent 工程。
  \item \textbf{产品经理与运营人员}：需要判断场景、设计体验、定义指标和验收系统。
  \item \textbf{独立开发者与创业者}：希望用有限资源做出真正可交付的 Agent 产品。
  \item \textbf{AI 学习者}：已经看过许多教程，但缺少一条从原理到项目的完整路径。
\end{itemize}

代码章节默认读者能看懂函数、字典、循环、异常处理和类型提示。你不需要先学完机器学习数学，也不需要训练自己的大模型。真正必要的前置知识是：能把一个问题拆成输入、过程、输出和验收条件。

\section*{怎样使用这本书}

第一遍按顺序读，每章运行示例并完成一个小任务。第二遍围绕毕业项目回查相关章节。不要试图同时学习所有框架；本书先用标准库构造最小实现，再说明何时引入 LangGraph、OpenAI Agents SDK、AutoGen、CrewAI、AgentScope 或其他框架。

书中有五类提示框：

\begin{description}[style=nextline,leftmargin=2em]
  \item[定义] 精确约束一个容易混淆的术语。
  \item[本节要点] 压缩必须保留的判断。
  \item[动手任务] 把理解转成可检查的产物。
  \item[风险提示] 标出权限、安全、成本或工程陷阱。
  \item[完成检查] 判断这一章是否真正学会。
\end{description}

\section*{关于模型、价格与时效}

Agent 生态变化很快。模型名称、API 形式、框架接口、Stars 和价格都会过期。因此正文尽量讲稳定机制，具体资源集中到附录并链接原始来源。版本相关信息以仓库更新和 GitHub Releases 为准。

\section*{版权与边界}

本书是原创导读和实践教程，不镜像第三方付费课程、电子书或视频。引用观点尽量回到论文、官方文档和作者原文。书稿采用 CC BY-NC-SA 4.0，示例代码和构建脚本采用 MIT；第三方项目仍受各自许可证约束。

\begin{practice}[开始前的承诺]
在一张纸或一个 Markdown 文件里写下：你希望 Agent 替谁完成什么任务？怎样算完成？哪一步出错最危险？把这三句话保留到全书结束，再看它们发生了什么变化。
\end{practice}

\begin{flushright}
DreamingRiverForUs\\
\BookDate
\end{flushright}
